%% start of file `template.tex'.
%% Copyright 2006-2013 Xavier Danaux (xdanaux@gmail.com).
%
% This work may be distributed and/or modified under the
% conditions of the LaTeX Project Public License version 1.3c,
% available at http://www.latex-project.org/lppl/.


\documentclass[11pt,a4paper]{moderncv}        % possible options include font size ('10pt', '11pt' and '12pt'), paper size ('a4paper', 'letterpaper', 'a5paper', 'legalpaper', 'executivepaper' and 'landscape') and font family ('sans' and 'roman')

% modern themes
\moderncvstyle{banking}                            % style options are 'casual' (default), 'classic', 'oldstyle' and 'banking'
\moderncvcolor{purple}                                % color options 'blue' (default), 'orange', 'green', 'red', 'purple', 'grey' and 'black'
%\renewcommand{\familydefault}{\sfdefault}         % to set the default font; use '\sfdefault' for the default sans serif font, '\rmdefault' for the default roman one, or any tex font name
\nopagenumbers{}                                  % uncomment to suppress automatic page numbering for CVs longer than one page

% character encoding
\usepackage{fontawesome5}
\usepackage{fontspec}
\usepackage{tabularx}
\usepackage{ragged2e}
% if you are not using xelatex ou lualatex, replace by the encoding you are using
% \usepackage{CJKutf8}                              % if you need to use CJK to typeset your resume in Chinese, Japanese or Korean
% \usepackage{xeCJK}
% \setCJKmainfont{IPAMincho}
% \setCJKsansfont{IPAGothic}
% \setCJKmonofont{IPAGothic}
% \PassOptionsToPackage{}{hyperref}
\usepackage[dvipsnames]{xcolor}
\definecolor{deepPurple}{RGB}{85,26,139}
\newcommand{\highlight}[1]{\textcolor{deepPurple}{\textbf{#1}}}

\usepackage{emoji}
\usepackage{luatexja-fontspec}

\setmainfont{Arial}                    % 设置英文无衬线字体
\setsansfont{Arial}                    % 如果想为无衬线设置单独字体，可以用这行
\setmonofont{Courier New}              % 英文等宽字体
% \setCJKmainfont{Noto Sans CJK JP}      % 设置中文主字体为无衬线字体
% \setCJKsansfont{Noto Sans CJK JP}      % 中文无衬线字体
% \setCJKmonofont{Noto Sans Mono CJK JP} % 中文等宽字体
\setmainjfont{Noto Sans CJK JP}        % 主体日文/中文
\setsansjfont{Noto Sans CJK JP}        % 无衬线
\setmonojfont{Noto Sans Mono CJK JP}   % 等宽

% adjust the page margins
\usepackage[scale=0.8]{geometry}
\usepackage{multicol}
%\setlength{\hintscolumnwidth}{3cm}                % if you want to change the width of the column with the dates
%\setlength{\makecvtitlenamewidth}{10cm}           % for the 'classic' style, if you want to force the width allocated to your name and avoid line breaks. be careful though, the length is normally calculated to avoid any overlap with your personal info; use this at your own typographical risks...

\usepackage{import}
\usepackage{todonotes}
\usepackage{enumitem}
\usepackage{fancyhdr}
\pagestyle{fancy}
\usepackage{xspace}

\newcommand{\strongit}[1]{\textbf{\textit{#1}}}

\fancyhf{}
\renewcommand{\headrulewidth}{0.4pt}       % 页眉细线
\fancyhead[L]{\textsc{\color{gray}[1] Curriculum Vitae with the Candidate's Photograph}} % 左页眉，小型大写+灰色
\fancyhead[R]{\textcolor{gray}{\thepage}}        % 右页眉显示页码


% personal data
\name{Deyun}{Lyu}
% \title{Curriculum Vitae}                               % optional, remove / comment the line if not wanted
% \address{}{}{}% optional, remove / comment the line if not wanted; the "postcode city" and and "country" arguments can be omitted or provided empty
% \phone[mobile]{909-839-3097}                   % optional, remove / comment the line if not wanted
% \phone[fixed]{01234 123456}                    % optional, remove / comment the line if not wanted
%\phone[fax]{+3~(456)~789~012}                      % optional, remove / comment the line if not wanted
% \email{xpan1@swarthmore.edu}                               % optional, remove / comment the line if not wanted
% \homepage{shawnpan.me}                         % optional, remove / comment the line if not wanted
% \extrainfo{}                 % optional, remove / comment the line if not wanted
\photo[64pt][0.4pt]{myPhoto}                       % optional, remove / comment the line if not wanted; '64pt' is the height the picture must be resized to, 0.4pt is the thickness of the frame around it (put it to 0pt for no frame) and 'picture' is the name of the picture file
% \quote{Some quote}                                 % optional, remove / comment the line if not wanted

% to show numerical labels in the bibliography (default is to show no labels); only useful if you make citations in your resume
%\makeatletter
%\renewcommand*{\bibliographyitemlabel}{\@biblabel{\arabic{enumiv}}}
%\makeatother
%\renewcommand*{\bibliographyitemlabel}{[\arabic{enumiv}]}% CONSIDER REPLACING THE ABOVE BY THIS

% bibliography with mutiple entries
%\usepackage{multibib}
%\newcites{book,misc}{{Books},{Others}}

\newcommand*{\customcventry}[7][.25em]{
  \begin{tabular}{@{}l} 
    {\bfseries #4}
  \end{tabular}
  \hfill% move it to the right
  \begin{tabular}{l@{}}
     {\bfseries #5}
  \end{tabular} \\
  \begin{tabular}{@{}l} 
    {\itshape #3}
  \end{tabular}
  \hfill% move it to the right
  \begin{tabular}{l@{}}
     {\itshape #2}
  \end{tabular}
  \ifx&#7&%
  \else{\\%
    \begin{minipage}{\maincolumnwidth}%
      \small#7%
    \end{minipage}}\fi%
  \par\addvspace{#1}}


\newcommand*{\customnewcventry}[8][.25em]{%
  \begin{tabular}{@{}l} 
    {\bfseries #4} % University
  \end{tabular}
  \hfill% move it to the right
  \begin{tabular}{l@{}}
     {\bfseries #5} % Date
  \end{tabular} \\
  \begin{tabular}{@{}l} 
    {\itshape #6} % Department
  \end{tabular}
  \hfill \\
  \begin{tabular}{@{}l} 
    {\itshape #3} % Degree/Position
  \end{tabular}
  \hfill% move it to the right
  \begin{tabular}{l@{}}
     {\itshape #2} % Location
  \end{tabular}
  \ifx&#8&%
  \else{\\%
    \begin{minipage}{\maincolumnwidth}%
      \small#8%
    \end{minipage}}\fi%
  \par\addvspace{#1}}

\newcommand*{\customcvproject}[4][.25em]{
%   \vfill\noindent
  \begin{tabular}{@{}l} 
    {\bfseries #2}
  \end{tabular}
  \hfill% move it to the right
  \begin{tabular}{l@{}}
     {\itshape #3}
  \end{tabular}
  \ifx&#4&%
  \else{\\%
    \begin{minipage}{\maincolumnwidth}%
      \small#4%
    \end{minipage}}\fi%
  \par\addvspace{#1}}

% % === Compact award entry ===
% \newcommand*\awardcventry[4]{%
%   % #1 = Date (e.g., "June 2025")
%   % #2 = Award/Grant title
%   % #3 = Organization
%   % #4 = (optional) small note, leave empty if none
%   \begin{tabularx}{\textwidth}{@{}X r@{}}
%     \textbf{#2}\\[-0.2ex]{\itshape #3} & {\itshape #1}
%   \end{tabularx}%
%   \ifx&#4&%
%   \else\\[-0.2ex]\begin{minipage}{\textwidth}\small #4\end{minipage}\fi
%   \par\addvspace{0.6em}%
% }

% ====== Awards (Minimal) ======
\usepackage{tabularx,xcolor}
\newcommand*\awardentry[4]{%
  % #1 Date   #2 Title   #3 Organization   #4 Optional note
  \noindent
  \begin{tabularx}{\textwidth}{@{}>{\bfseries}X r@{}}
    #2 & {\itshape #1} \\
  \end{tabularx}\\[-0.2ex]
  {\itshape #3}%
  \if\relax\detokenize{#4}\relax\else\\[-0.25ex]{\footnotesize\color{gray!65}#4}\fi
  \par\addvspace{0.75em}%
}

% ==============================
% 学生指导条目（自动带圆点）
% 用法：\supervisionentry{学生名}{时间}{学校}{地点}{论文题目}
% ==============================
\newcommand{\supervisionentry}[5]{%
  \item \textbf{#1}, #2, #3, #4 \\ % 第一行
  {#5} % 第二行（论文题目斜体）
}

\setlength{\tabcolsep}{12pt}

% % 这是 LaTeX 的一个特殊命令，允许访问和修改以 @ 开头的内部宏（这里就是 \@firstname、\@familyname 等）
% \makeatletter
% % 重新定义 moderncv 模板中的 \makecvtitle 命令
% \renewcommand*{\makecvtitle}{%
%   % 让标题整体 居中
%   \begin{center}
%     % ----- name (use theme color) -----
%     {\textcolor{color1}{\Huge\bfseries\@firstname~\@familyname}}\par
%     % spacing name <-> address
%     \vspace*{4mm}
%     % ----- address -----
%     \ifx\@addressstreet\@empty\else
%       {\large\@addressstreet}\par
%     \fi
%     % spacing address <-> next content
%     \vspace*{10mm}
%   \end{center}
% }
% \makeatother

% \makeatletter
% \renewcommand*{\makecvtitle}{%
%   \begin{center}
%     % ----- name -----
%     {\textcolor{color1}{\Huge\bfseries\@familyname~\@firstname}}\par

%     % ----- title / position -----
%     \vspace*{2mm}
%     {\normalsize Research Fellow | National Institute of Informatics}\par

%     % ----- address -----
%     \vspace*{2mm}
%     {\normalsize 2-1-2 Hitotsubashi, Chiyoda-ku, Tokyo, Japan}\par

%     % ----- contact info -----
%     \vspace*{2mm}
%     \begin{tabular}{c c c c}
%       \faPhone\enspace +81 080-2567-8231 &
%       \faEnvelope \enspace lyudeyun@nii.ac.jp & 
%       \faLink\enspace lyudeyun.github.io & 
%       % \faGithub\enspace lyudeyun 
%     \end{tabular}
%   \end{center}
% }
% \makeatother


\makeatletter
\renewcommand*{\makecvtitle}{%
  \begin{tabularx}{\textwidth}{@{}X r@{}}
    % 左边：名字 + 联系方式
    \begin{minipage}[t]{0.80\linewidth}
      % ----- name -----
      {\textcolor{color1}{\Huge\bfseries\@familyname~\@firstname}}\par

      % ----- title / position -----
      \vspace*{2mm}
      {\normalsize Research Fellow | National Institute of Informatics}\par

      % ----- nationality and birth -----
      \vspace*{2mm}
      \begin{tabular}{{@{}lll}}
        \faGlobeAsia \enspace Nationality: China &
        \faBirthdayCake \enspace Date of birth: 1997/01/23
      \end{tabular}

      % ----- contact info -----
      % \vspace*{2mm}
      % \begin{tabular}{{@{}lll}}
      %   \faPhone\enspace +81 080-2567-8231 & \faEnvelope\enspace lyudeyun@nii.ac.jp &
      %   \faHome\enspace lyudeyun.github.io
      % \end{tabular}
      \vspace*{2mm}
      \begin{tabular}{@{}l@{\hspace{8pt}}l@{\hspace{8pt}}l@{}}
      \faPhone\enspace 080-2567-8231 & 
      \faEnvelope\enspace lyudeyun@nii.ac.jp & 
      \faHome\enspace lyudeyun.github.io
      \end{tabular}
      
      % ----- address -----
      \vspace*{2mm}
      {\normalsize \faLocationArrow\enspace 2-1-2 Hitotsubashi, Chiyoda-ku, Tokyo 101-8430, Japan}\par
    \end{minipage}
    &
    % 右边：照片，垂直居中
    \raisebox{-0.7\height}{\includegraphics[height=120pt]{myphoto.jpg}}
  \end{tabularx}
}
\makeatother

% 导言区
\makeatletter
\newcommand{\subsubsection}[1]{%
  \par\vspace{0.5em}%
  {\color{color1}\bfseries #1}%
  \par\vspace{0.25em}%
}
\makeatother


%----------------------------------------------------------------------------------
%            content
%----------------------------------------------------------------------------------
\begin{document}
%\begin{CJK*}{UTF8}{gbsn}                          % to typeset your resume in Chinese using CJK
%-----       resume       ---------------------------------------------------------
\makecvtitle
\vspace*{-5mm}

\section{About Me}

I am a Research Fellow at the \href{https://www.nii.ac.jp}{National Institute of Informatics}, working in the Trustworthy $\&$ Smart Software Engineering Lab led by Prof. \href{https://research.nii.ac.jp/~f-ishikawa}{\textbf{Fuyuki Ishikawa}}. I received my Ph.D. in Information Science and Electrical Engineering from Kyushu University under the supervision of Prof. \href{https://stap.ait.kyushu-u.ac.jp/~zhao/}{\textbf{Jianjun Zhao}}, with co-advisors Prof. \href{https://group-mmm.org/~arcaini/}{\textbf{Paolo Arcaini}} and Prof. \href{https://choshina.github.io/}{\textbf{Zhenya Zhang}}. My research interests focus on the quality assurance of AI-enabled cyber-physical systems (AI-CPSs).

\section{Education}

{\customcventry{Sept. 2011 -- June 2014}{High School Diploma in Science Program}{Hebei Wuyi Senior High School}{Hengshui, China}{}{\hfill \textit{Graduated}}}
\vspace{1mm}
{\customcventry{Sept. 2014 -- June 2018}{B.S. in Network Engineering}{Dalian University of Technology}{Dalian, China}{}{Supervisor: \href{https://faculty.dlut.edu.cn/gaojing/en/zsxx/960677/list/index.htm}{Prof. Jing Gao}\hfill \textit{Graduated}}}
\vspace{1mm}
{\customcventry{Sept. 2018 -- June 2020}{M.S. in Software Engineering}{Dalian University of Technology}{Dalian, China}{}{Supervisor: \href{https://faculty.dlut.edu.cn/kongweiqiang/en/more/828473/gzjlgd/index.htm}{Prof. Weiqiang Kong} \hfill \textit{Graduated}}}
\vspace{1mm}
{\customcventry{Oct. 2020 -- Mar. 2025}{Ph.D. in Advanced Information Technology}{Kyushu University}{Fukuoka, Japan}{}{Supervisor: \href{https://stap.ait.kyushu-u.ac.jp/~zhao/}{Prof. Jianjun Zhao}, Co-advisors: \href{https://group-mmm.org/~arcaini/}{Prof. Paolo Arcaini}, \href{https://choshina.github.io/}{Prof. Zhenya Zhang} \hfill \textit{Graduated}}}


% 导言区（可选，控制间距）
\setlist[itemize]{leftmargin=0em,itemsep=0.15em,topsep=0.2em}
% 正文
\section{Academic Degree}
\begin{itemize}[label={}]
    \item \textbf{Bachelor of Engineering} \hfill \textit{June 2018}\\
        School of Software Technology, Dalian University of Technology\\
        Major: Network Engineering \\
        Thesis: \emph{Design and Implementation of a Campus Second-hand Trading System} 
    \item \textbf{Master of Software Engineering} \hfill \textit{June 2020}\\
        School of Software Technology, Dalian University of Technology\\
        Major: Software Engineering \\
        Thesis: \emph{Research and Implementation of Infrared Binocular Camera Calibration Method} 
    \item \textbf{Doctor of Information Science} \hfill \textit{Mar. 2025}\\
        Graduate School of Information Science and Electrical Engineering, Kyushu University \\
        Major: Advanced Information Technology \\
        Dissertation: \emph{Quality Assurance of AI-Enabled Cyber-Physical Systems} \\ 
        (Selected as IPSJ SIG-SE Recommended Doctoral Dissertation \emoji{trophy}) 
\end{itemize}

\section{Work Experience}

{\customnewcventry{Apr. 2025 -- Present}{Research Fellow}{National Institute of Informatics}{Tokyo, Japan}{Information Systems Architecture Science Research Division}{}{}}
{\customnewcventry{Summer 2024}{Teaching Assistant}{Kyushu University}{Fukuoka, Japan}{Graduate School of Information Science and Electrical Engineering}{}{Course: \textit{Python Programming Exercises (Pythonプログラミング演習)} for undergraduate students}}
{\customnewcventry{Sept. 2023 -- Mar. 2024}{Visiting Researcher}{National Institute of Informatics}{Tokyo, Japan}{Information Systems Architecture Science Research Division}{}{}}
{\customnewcventry{Summer 2023}{Teaching Assistant}{Kyushu University}{Fukuoka, Japan}{Graduate School of Information Science and Electrical Engineering}{}{Course: \textit{Python Programming Exercises (Pythonプログラミング演習)} for undergraduate students}}
{\customcventry{Oct. 2021 -- Sept. 2023}{Research Fellow (JST-SPRING)}{Japan Science and Technology Agency}{Fukuoka, Japan}{}{}}
% {\customcventry{Summer 2021}{Teaching Assistant}{Kyushu University}{Fukuoka, Japan}{Graduate School of Information Science and Electrical Engineering}{}{Course: \textit{Machine Learning Systems Engineering (機械学習システム工学)} for graduate students}}
{\customnewcventry{Oct. 2020 -- Sept. 2021}{Research Assistant}{Kyushu University}{Fukuoka, Japan}{Graduate School of Information Science and Electrical Engineering}{}{}}
{\customnewcventry{Fall 2019}{Teaching Assistant}{Dalian University of Technology}{Dalian, China}{School of Software Technology}{}{Course: \textit{C++ Programming Exercises} for undergraduate students}}
{\customnewcventry{Spring 2019}{Teaching Assistant}{Dalian University of Technology}{Dalian, China}{School of Software Technology}{}{Course: \textit{C Programming Exercises} for undergraduate students}}
{\customnewcventry{Fall 2018}{Teaching Assistant}{Dalian University of Technology}{Dalian, China}{School of Software Technology}{}{Course: \textit{Operating Systems Practicum} for undergraduate students}}
{\customnewcventry{Sept. 2018 -- June. 2020}{Research Assistant}{Dalian University of Technology}{Dalian, China}{School of Software Technology}{}{}}

\setlist[itemize]{leftmargin=1.2em,itemsep=0.15em,topsep=0.2em}

\section{Specialized Fields}

\begin{itemize}
  \item Software Engineering
  \item Formal Methods
  \item Cyber-Physical Systems
  \item Artificial Intelligence 
\end{itemize}

\subsection{Research Interests}
\begin{itemize}
  \item Trustworthy AI-Enabled Cyber-Physical Systems
  \item Safe and Reliable Autonomous Driving Systems
  \item Software Engineering for AI Systems (SE4AI)
  \item Formal Verification and Validation
\end{itemize}

\section{Professional Memberships and Services}

\subsection{Academic Memberships}
\begin{itemize}
  % \item Member, ACM (Association for Computing Machinery) 
  \item Member, IEEE (Institute of Electrical and Electronics Engineers) \hfill  
  \item Member, IPSJ (Information Processing Society of Japan)
  \item Member, IPSJ SIG-SE (Special Interest Group on Software Engineering)
\end{itemize}

\subsection{Professional Services}

\subsubsection{Journal Paper}

\begin{itemize}
    \item Reviewer, ACM Transactions on Software Engineering and Methodology (TOSEM)
    % \item Reviewer, IEEE Transactions on Software Engineering (TSE)
    \item Reviewer, Empirical Software Engineering (EMSE)
    \item Reviewer, IEEE Transactions on Reliability (TRel)
    \item Reviewer, Research Directions: Cyber-Physical Systems
\end{itemize}

\subsubsection{Conference Paper}

\begin{itemize}
    \item Web Chair, \href{https://conf.researchr.org/committee/fm-2026/fm-2026-organizing-committee}{International Symposium on Formal Methods (FM 2026)}
    \item PC Member, \href{https://conf.researchr.org/committee/atva-2025/atva-2025-artifact-evaluation-committee}{International Symposium on Automated Technology for Verification and Analysis (ATVA 2025) – Artifact Evaluation Track}
    \item PC Member, Association for the Advancement of Artificial Intelligence (AAAI 2023, 2024, 2025)
    \item Reviewer, International Conference on Software Engineering (ICSE)
    \item Reviewer, IEEE/ACM International Conference on Automated Software Engineering (ASE)
    \item Reviewer, ACM SIGSOFT Symposium on the Foundations of Software Engineering (FSE)
    \item Reviewer, International Symposium on Software Testing and Analysis (ISSTA)
    \item Reviewer, International Conference on Machine Learning (ICML)
    \item Reviewer, International Conference on Embedded Software (EMSOFT)
    \item Reviewer, IEEE International Conference on Software Testing, Verification and Validation (ICST)
    \item Reviewer, IEEE/RSJ International Conference on Intelligent Robots and Systems (IROS)    
    \item Reviewer, International Conference on Engineering of Complex Computer Systems (ICECCS)
    \item Reviewer, International Symposium on Theoretical Aspects of Software Engineering (TASE)
    \item Reviewer, International Conference on AI Engineering – Software Engineering for AI (CAIN)
\end{itemize}



\section{History of Awards and Disciplinary Actions}

\subsection{Awards}

\awardentry{Aug. 2025}
{2025 IPSJ SIG-SE Excellent Research Award (情報処理学会 ソフトウェア工学研究会 2025年度卓越研究賞)}
{Information Processing Society of Japan, Special Interest Group on Software Engineering (IPSJ SIG-SE)}
{For our TOSEM 2025 paper entitled ``SpectAcle: Fault Localisation of AI-Enabled CPS by Exploiting Sequences of DNN Controller Inferences."}

\awardentry{June 2025}
{2025 IPSJ SIG-SE Recommended Doctoral Dissertation（情報処理学会 ソフトウェア工学研究会 2025 年度推薦博士論文）}
{Information Processing Society of Japan, Special Interest Group on Software Engineering (IPSJ SIG-SE)}
{}

\awardentry{May 2024}
{GECCO 2024 Student Travel Grant}
{ACM Special Interest Group on Genetic and Evolutionary Computation (ACM SIGEVO)}
{(Amount: \$600 USD)}

\awardentry{Oct. 2021 -- Sept. 2023}
{Support for Pioneering Research Initiated by the Next Generation (JST SPRING, 次世代研究者挑戦的研究プログラム)}
{Japan Science and Technology Agency (JST)}
{(Grant No. JPMJSP2136; Amount: JPY 2,400,000 per year)}

\awardentry{Sept. 2021}
{Graduate School Research Support Scholarship}
{Kyushu University}
{(Amount: JPY 250,000)}

\awardentry{Sept. 2018 -- June 2020}
{First-Class Scholarship of the Graduate School}
{Dalian University of Technology}
{(Amount: RMB 19,600 per year)}

\subsection{Disciplinary Actions}

None

\clearpage
\fancyhf{}
\renewcommand{\headrulewidth}{0.4pt}       % 页眉细线
\fancyhead[L]{\textsc{\color{gray}[2] Achievements in Education}} % 左页眉，小型大写+灰色
\fancyhead[R]{\textcolor{gray}{\thepage}}        % 右页眉显示页码

\section{Achievements in Education}

During my doctoral and postdoctoral studies, I have gained teaching experience primarily through teaching assistantships and mentoring roles, assisting around \textbf{400} undergraduate and graduate students. These experiences encompassed managing large-capacity undergraduate courses and supervising graduate students in thesis projects related to machine learning and cyber-physical systems. I consistently emphasized project-based and hands-on learning, while fostering students' independent thinking, which helped them better connect theoretical concepts with real-world applications.

\subsection{Teaching Assistantships}

I served as a teaching assistant for undergraduate courses, where I conducted laboratory sessions, supported students with programming exercises, and prepared and graded assignments. These experiences strengthened my ability to guide students, provide constructive feedback, and address technical questions. Courses included:

\begin{itemize}
    \supervisionentry{Python Programming Exercises}{Summer 2024}{Kyushu University}{Fukuoka, Japan}{(approx. 60 undergraduate students)}
    \supervisionentry{Python Programming Exercises}{Summer 2023}{Kyushu University}{Fukuoka, Japan}{(approx. 30 undergraduate students)}
    % \supervisionentry{Machine Learning Systems Engineering}{Summer 2021}{Kyushu University}{Fukuoka, Japan}{(approx. 30 undergraduate students)}
    \supervisionentry{C++ Programming Exercises}{Fall 2019}{Dalian University of Technology}{Dalian, China}{(approx. 100 undergraduate students)}
    \supervisionentry{C Programming Exercises}{Spring 2019}{Dalian University of Technology}{Dalian, China}{(approx. 100 undergraduate students)}
    \supervisionentry{Operating Systems Practicum}{Fall 2018}{Dalian University of Technology}{Dalian, China}{(approx. 100 undergraduate students)}
\end{itemize}

A major challenge in undergraduate programming courses was that many students struggled with understanding abstract concepts in Python and C/C++. Common difficulties included not grasping the difference between value and reference, repeatedly making the same syntax errors, or failing to design programs that went beyond simple “copy-and-paste” solutions. To address this, I designed laboratory sessions that focused on incremental problem-solving: \strongit{breaking larger assignments into smaller, testable components}. I also \strongit{encouraged students to engage in pair programming}, which helped weaker students learn by observing stronger peers and motivated stronger students to consolidate their understanding by explaining concepts. Weekly Q\&A sessions and office hours provided additional support, where I emphasized debugging as a learning process rather than a source of frustration.

\subsection{Mentoring Experience}

I have supervised three graduate students who successfully completed their theses in machine learning and cyber-physical systems, and I am currently supervising two additional graduate students. I pay special attention to fostering their independent thinking and problem-solving abilities. In thesis supervision, I found that graduate students often had difficulty framing research questions or connecting their technical implementations to broader concepts in machine learning and CPS. This was partly because they focused too much on implementation details while neglecting higher-level abstraction and conceptualization of problems and solutions. To address this, I first encouraged them to articulate research questions and construct abstract models before diving into implementation. I then \strongit{encouraged them to decompose these problems into manageable hypotheses} and to \strongit{formulate precise research objectives}, enabling step-by-step validation of their ideas and methods. This not only strengthened their ability to think abstractly but also improved their technical rigour, while nurturing independent thinking and resilience in facing open-ended research challenges. Several of these projects were completed successfully as master’s theses, and some students plan to continue research in academia.

\begin{itemize}
    \supervisionentry{Chenchang Wang}{Expected to graduate in March 2027}{Kyushu University}{Fukuoka, Japan}{Master's thesis: \emph{eBPF-assisted Falsification of Temporal Logic Specifications under Noisy Environments}}
    \supervisionentry{Borui Zhang}{Expected to graduate in March 2026}{Kyushu University}{Fukuoka, Japan}{Master's thesis: \emph{ContrRep: Search-Based Repair of DNN Controllers of AI-Enabled CPSs}}
    \supervisionentry{Yipei Yan}{Graduated in March 2025}{Kyushu University}{Fukuoka, Japan}{Master's thesis: \emph{STL公式に基づくCPSテスト用ベンチマークの自動生成技術}}
    \supervisionentry{Yi Li}{Graduated in March 2025}{Kyushu University}{Fukuoka, Japan}{Master's thesis: \emph{Tactical: Fault Localization of AI-Enabled Cyber-Physical Systems by Exploiting Temporal Neuron Activation}}
    \supervisionentry{Ruinan Zhang}{Graduated in June 2022}{Dalian University of Technology}{Dalian, China}{Master's thesis: \emph{Research on Test Prioritization Technique for Deep Neural Networks}}
\end{itemize}

\clearpage

\fancyhf{}
\renewcommand{\headrulewidth}{0.4pt}       % 页眉细线
\fancyhead[L]{\textsc{\color{gray}[3] Achievements in Research}} % 左页眉，小型大写+灰色
\fancyhead[R]{\textcolor{gray}{\thepage}}        % 右页眉显示页码

\section{Refereed Publications}

\subsection{Quality Assurance of AI-Enabled Cyber-Physical Systems}

\begin{itemize}
    \item \hypertarget{pub:AI-CPS-P7}{}\textbf{[AI-CPS-P7]} \textbf{Deyun Lyu}, Jiayang Song, Zhenya Zhang, Zhijie Wang, Tianyi Zhang, Lei Ma and Jianjun Zhao. ``Repair of DNN Control Policies via STL-Guided Patch Synthesis." IEEE Transactions on Reliability. (TREL 2025, the flagship journal in reliability engineering, \highlight{CORE Rank A}, Impact Factor: 5.7, under review)
    \item \hypertarget{pub:AI-CPS-P6}{}\textbf{[AI-CPS-P6]} \textbf{Deyun Lyu}, Zhenya Zhang, Paolo Arcaini, Thomas Laurent, Borui Zhang, Fuyuki Ishikawa, and Jianjun Zhao. ``ContrRep: Search-Based Repair of DNN Controllers of AI-Enabled CPSs." Automated Software Engineering. (ASEJ 2025, \highlight{CORE Rank B}, Impact Factor: 3.1, under review)
    \item \hypertarget{pub:AI-CPS-P5}{}\textbf{[AI-CPS-P5]} \textbf{Deyun Lyu}, Yi Li, Zhenya Zhang, Paolo Arcaini, Xiao-Yi Zhang, Fuyuki Ishikawa and Jianjun Zhao. ``\href{https://www.sciencedirect.com/science/article/abs/pii/S0164121225001438}{Fault Localisation of AI-Enabled Cyber-Physical Systems by Exploiting Temporal Neuron Activation.}" Journal of Systems and Software. (JSS 2025, \highlight{CORE Rank A}, Impact Factor: 4.1) 
    \item \hypertarget{pub:AI-CPS-P4}{}\textbf{[AI-CPS-P4]} \textbf{Deyun Lyu}, Zhenya Zhang, Paolo Arcaini, Xiao-Yi Zhang, Fuyuki Ishikawa and Jianjun Zhao. ``\href{https://dl.acm.org/doi/10.1145/3705307}{SpectAcle: Fault Localisation of AI-Enabled CPS by Exploiting Sequences of DNN Controller Inferences.}" ACM Transactions on Software Engineering and Methodology. (TOSEM 2025, a top-tier journal in software engineering, \highlight{CORE Rank A*}, Impact Factor: 6.6, \highlight{IPSJ SIG-SE Excellent Research Award})
    \item \hypertarget{pub:AI-CPS-P3}{}\textbf{[AI-CPS-P3]} \textbf{Deyun Lyu}, Zhenya Zhang, Paolo Arcaini, Fuyuki Ishikawa, Thomas Laurent and Jianjun Zhao. ``\href{https://dl.acm.org/doi/10.1145/3638529.3654078}{Search-Based Repair of DNN Controllers of AI-Enabled CPS Guided by System-Level Specifications.}"  Proceedings of the Genetic and Evolutionary Computation Conference. (GECCO 2024, a top-tier conference in computational intelligence, \highlight{CORE Rank A})
    \item \hypertarget{pub:AI-CPS-P2}{}\textbf{[AI-CPS-P2]} Zhenya Zhang, \textbf{Deyun Lyu}, Paolo Arcaini, Lei Ma, Ichiro Hasuo and Jianjun Zhao. ``\href{https://ieeexplore.ieee.org/document/9844247}{FalsifAI: Falsification of AI-Enabled Hybrid Control Systems Guided by Time-Aware Coverage Criteria.}" IEEE Transactions on Software Engineering. (TSE 2023, the flagship journal in software engineering, \highlight{CORE Rank A*}, Impact Factor: 6.5)
    \item \hypertarget{pub:AI-CPS-P1}{}\textbf{[AI-CPS-P1]} Jiayang Song, \textbf{Deyun Lyu}*, Zhenya Zhang, Zhijie Wang, Tianyi Zhang and Lei Ma. ``\href{https://dl.acm.org/doi/10.1145/3510457.3513049}{When Cyber-Physical Systems Meet AI: A Benchmark, an Evaluation, and a Way Forward.}" 44th International Conference on Software Engineering: Software Engineering in Practice. (ICSE 2022, the flagship conference in software engineering, \highlight{Co-first author}, \highlight{CORE Rank A*})
\end{itemize}

\subsection{Quality Assurance of Cyber-Physical Systems}

\begin{itemize}
    \item \textbf{[CPS-P4]} Tanmay Khandait, \textbf{Deyun Lyu}*, (12 other authors). ``ARCH-COMP 2025 Category Report: Falsification.`` 12th International Workshop on Applied Verification of Continuous and Hybrid Systems. (\highlight{Co-first author}, to appear)
    \item \hypertarget{pub:CPS-P3}{}\textbf{[CPS-P3]} Yipei Yan, \textbf{Deyun Lyu}, Zhenya Zhang, Paolo Arcaini and Jianjun Zhao. ``\href{https://ieeexplore.ieee.org/document/10922764}{Automated Generation of Benchmarks for Falsification of STL Specifications.}" IEEE Transactions on Computer-Aided Design of Integrated Circuits and Systems. (TCAD 2025, the flagship journal in electronic design automation, \highlight{CCF A}, Impact Factor: 2.9)
    \item \hypertarget{pub:CPS-P2}{}\textbf{[CPS-P2]} Zhenya Zhang, \textbf{Deyun Lyu}, Paolo Arcaini, Lei Ma, Ichiro Hasuo and Jianjun Zhao. ``\href{https://link.springer.com/chapter/10.1007/978-3-030-81685-8_29}{Effective Hybrid System Falsification Using Monte Carlo Tree Search Guided by QB-Robustness.}" 33rd International Conference on Computer-Aided Verification. (CAV 2021, the flagship conference in formal methods, \highlight{CORE Rank A*})
    \item \textbf{[CPS-P1]} Zhenya Zhang, \textbf{Deyun Lyu}, Paolo Arcaini, Lei Ma, Ichiro Hasuo and Jianjun Zhao. ``\href{https://link.springer.com/chapter/10.1007/978-3-030-76384-8_24}{On the Effectiveness of Signal Rescaling in Hybrid System Falsification.}" 13th NASA Formal Methods Symposium. (NFM 2021)
\end{itemize}

\subsection{Quality Assurance of Autonomous Driving Systems} 
\noindent \textbf{Note:} \textit{These works were conducted during my time at Dalian University of Technology. While my official academic name is ``Deyun Lyu", one of them ([ADS-Verf-P1]) lists me as ``Deyun Lv" following pinyin conventions.}

\begin{itemize}
    \item \textbf{[ADS-Verf-P2]} Huihui Wu, \textbf{Deyun Lyu}, Y. Zhang, Gang Hou, Masahiko Watanabe, Jie Wang and Weiqiang Kong. ``\href{https://ietresearch.onlinelibrary.wiley.com/doi/full/10.1049/itr2.12162}{A Verification Framework for Behavioural Safety of Self-Driving Cars.}" IET Intelligent Transport Systems. (\highlight{CCF C}, Impact Factor: 2.3)
    \item \textbf{[ADS-Verf-P1]} Huihui Wu, \textbf{Deyun Lyu}, Tengxiang Cui, Gang Hou, Masahiko Watanabe and Weiqiang Kong. ``\href{https://link.springer.com/article/10.1007/s00165-021-00539-2}{SDLV: Verification of Steering Angle Safety for Self-Driving Cars.}" Formal Aspects of Computing. (FAC 2021, \highlight{CCF B}, Impact Factor: 1.8)
\end{itemize}

\section{Invited Talks, Presentations, and Posters}

\subsection{Invited Talks}

\begin{itemize}
    \item \textbf{[IT3]} ``Quality Assurance of AI-Enabled Cyber-Physical Systems." The 220th IPSJ/SIGSE Software Engineering Workshop (第221回 ソフトウェア工学研究発表会), 東大寺総合文化センター, Nara, Japan. (\textit{Scheduled, Nov. 2025})
    % \item \textbf{[IT3]} ``Quality Assurance of AI-Enabled Cyber-Physical Systems." IPSJ/SIGSE Software Engineering Symposium (ソフトウェアエンジニアリングシンポジウム 2025), Waseda University, Tokyo, Japan. (\textit{Sept. 2025})
    \item \textbf{[IT2]} ``SpectAcle: Fault Localisation of AI-Enabled CPS by Exploiting Sequences of DNN Controller Inferences." \href{https://ses.sigse.jp/2025/program.html}{IPSJ/SIGSE Software Engineering Symposium (ソフトウェアエンジニアリングシンポジウム 2025)}, Waseda University, Tokyo, Japan. (\textit{Sept. 2025})
    \item \textbf{[IT1]} ``Search-Based Repair of DNN Controllers of AI-Enabled Cyber-Physical Systems Guided by System-Level Specifications." Top Conference/Journal Special Talk (トップカンファレンス・トップ論文誌特別講演), \href{https://jssst2024.wordpress.com/program/}{41st Annual Conference of the Japan Society for Software Science and Technology (日本ソフトウェア科学会第41回大会)}, Ritsumeikan University, Osaka, Japan. (\textit{Sept. 2024})
\end{itemize}

\subsection{Presentations}

\begin{itemize}
    \item \textbf{[Pres4]} ``SpectAcle: Fault Localisation of AI-Enabled CPS by Exploiting Sequences of DNN Controller Inferences." Journal-First Track, 40th IEEE/ACM International Conference on Automated Software Engineering (ASE 2025), Seoul, South Korea. (\textit{Scheduled, Nov. 2025})
    \item \textbf{[Pres3]} ``Fault Localisation of AI-Enabled Cyber-Physical Systems by Exploiting Temporal Neuron Activation." Journal-First Track, 40th IEEE/ACM International Conference on Automated Software Engineering (ASE 2025), Seoul, South Korea. (\textit{Scheduled, Nov. 2025})
    \item \textbf{[Pres2]} ``Search-Based Repair of DNN Controllers of AI-Enabled CPS Guided by System-Level Specifications." Proceedings of the Genetic and Evolutionary Computation Conference (GECCO 2024), Melbourne, Australia. (\textit{July 2025})
    \item \textbf{[Pres1]} ``Towards Building Reliable AI-Enabled Cyber-Physical Systems." Doctoral Symposium, 26th International Symposium on Formal Methods (FM 2023), Lübeck, Germany. (\textit{Mar. 2023, Accepted, but not presented due to COVID-19 travel restrictions.})
\end{itemize}

\subsection{Posters}

\begin{itemize}
    \item \textbf{[Poster1]} ``FalsifAI: Falsification of AI-Enabled Hybrid Control Systems Guided by Time-Aware Coverage Criteria." Demo/Poster Session (デモ・ポスターセッション), 39th Annual Conference of the Japan Society for Software Science and Technology (日本ソフトウェア科学会第39回大会), Nanzan University, Nagoya, Japan. (\textit{Sept. 2022})
\end{itemize}

\section{Books and Book Chapters}
None

\section{Five Representative Publications}


My doctoral research focuses on the \textbf{quality assurance of AI-enabled cyber-physical systems} (AI-CPSs), a critical challenge for ensuring the reliability and trustworthiness of AI controllers embedded in physical systems such as autonomous vehicles and robots. These five publications collectively span the lifecycle of quality assurance of such systems, including benchmarking [AI-CPS-P1], testing [AI-CPS-P2], fault localisation [AI-CPS-P4, AI-CPS-P5], and repair [AI-CPS-P3]. Prior to this, existing research often missed the forest for the trees (木を見て、森を見ず): (i) no systematic benchmarking or evaluation of AI-CPSs existed; (ii) testing focused only on system-level correctness while ignoring DNN behaviour; and (iii) fault localisation and repair targeted simple DNN classifiers, failing to capture system-level complexity. In contrast, my research \textbf{sees both the trees and the forest} (木も見て、森も見る), establishing a \textbf{system–component co-analysis} methodology that advances quality assurance and paves the way for trustworthy AI-CPSs. Beyond academic contributions, these techniques are being advanced into practice through the collaboration with Toyota and the Japan Automobile Research Institute, toward deployment in their end-to-end autonomous driving systems.

\begin{itemize}
    \item \textbf{[AI-CPS-P1]} Jiayang Song, \textbf{Deyun Lyu}*, Zhenya Zhang, Zhijie Wang, Tianyi Zhang and Lei Ma. ``\href{https://dl.acm.org/doi/10.1145/3510457.3513049}{When Cyber-Physical Systems Meet AI: A Benchmark, an Evaluation, and a Way Forward.}" 44th International Conference on Software Engineering: Software Engineering in Practice. (ICSE 2022, the flagship conference in software engineering, \highlight{Co-first author}, \highlight{CORE Rank A*})
    
    \textbf{Summary}: Presents \strongit{nine pioneering industrial benchmarks} of AI-CPSs across seven domains, \strongit{filling the long-standing gap of systematic evaluation}. Enables multi-view comparisons between AI and classic controllers, showing hybrid controllers superior and motivating AI-aware testing. This work \strongit{laid the benchmarking foundation} for subsequent trustworthy AI-CPS research.
    \item \textbf{[AI-CPS-P2]} Zhenya Zhang, \textbf{Deyun Lyu}, Paolo Arcaini, Lei Ma, Ichiro Hasuo and Jianjun Zhao. ``\href{https://ieeexplore.ieee.org/document/9844247}{FalsifAI: Falsification of AI-Enabled Hybrid Control Systems Guided by Time-Aware Coverage Criteria.}" IEEE Transactions on Software Engineering. (TSE 2023, the flagship journal in software engineering, \highlight{CORE Rank A*}, Impact Factor: 6.5)

    \textbf{Summary}: Proposes \strongit{eight time-aware coverage criteria} tailored for AI-CPS and FalsifAI, the \strongit{first} coverage-guided falsification framework balancing \strongit{exploration} and \strongit{exploitation}. Unlike exploitation-only falsification approaches, it systematically uncovers diverse temporal behaviours. It significantly improves falsification effectiveness and \strongit{pioneered coverage-guided testing for AI-CPSs}, inspiring later testing research.
    \item \textbf{[AI-CPS-P4]} \textbf{Deyun Lyu}, Zhenya Zhang, Paolo Arcaini, Xiao-Yi Zhang, Fuyuki Ishikawa and Jianjun Zhao. ``\href{https://dl.acm.org/doi/10.1145/3705307}{SpectAcle: Fault Localisation of AI-Enabled CPS by Exploiting Sequences of DNN Controller Inferences.}" ACM Transactions on Software Engineering and Methodology. (TOSEM 2025, a top-tier journal in software engineering, \highlight{CORE Rank A*}, Impact Factor: 6.6, \highlight{IPSJ SIG-SE Excellent Research Award})
    
    \textbf{Summary}: Introduces the \strongit{first} fault localisation approach for AI-CPSs. Addressing the lack of per-inference ground truth in sequential control decisions, it exploits temporal sequences of DNN inferences with suspiciousness metrics to \strongit{identify faulty weights}. Evaluated on \strongit{6,067} faulty benchmarks, it effectively localises faults, \strongit{enabling subsequent analysis and repair of AI-CPSs}.
    \item \textbf{[AI-CPS-P5]} \textbf{Deyun Lyu}, Yi Li, Zhenya Zhang, Paolo Arcaini, Xiao-Yi Zhang, Fuyuki Ishikawa and Jianjun Zhao. ``\href{https://www.sciencedirect.com/science/article/abs/pii/S0164121225001438}{Fault Localisation of AI-Enabled Cyber-Physical Systems by Exploiting Temporal Neuron Activation.}" Journal of Systems and Software. (JSS 2025, \highlight{CORE Rank A}, Impact Factor: 4.1) 

    \textbf{Summary}: Proposes a neuron-level fault localisation approach for AI-CPSs. By linking temporal neuron activations with system-level specifications, it overcomes the absence of ground truth and identifies faulty neurons. It complements weight-level approaches, \strongit{offering an additional scale of analysis within fault localisation for AI-CPSs}.
    \item \textbf{[AI-CPS-P3]} \textbf{Deyun Lyu}, Zhenya Zhang, Paolo Arcaini, Fuyuki Ishikawa, Thomas Laurent and Jianjun Zhao. ``\href{https://dl.acm.org/doi/10.1145/3638529.3654078}{Search-Based Repair of DNN Controllers of AI-Enabled CPS Guided by System-Level Specifications.}"  Proceedings of the Genetic and Evolutionary Computation Conference. (GECCO 2024, a top-tier conference in computational intelligence, \highlight{CORE Rank A})

    \textbf{Summary}: Introduces the \strongit{first} search-based repair approach for DNN controllers of AI-CPSs guided by system-level specifications. Addressing the lack of ground truth, it searches for alternative values of suspicious weights identified through fault localisation, and demonstrates effective, efficient repair that \strongit{strengthens the reliability of AI-CPSs}.
\end{itemize}

\section{Research Grants and Funding}

\begin{itemize}
    \item Principal Investigator, Support for Pioneering Research Initiated by the Next Generation (次世代研究者挑戦的研究プログラム), Japan Science and Technology Agency (JST), Oct. 2021 -- Sept. 2023. Grant No. JPMJSP2136. Amount: JPY 500,000 annually (total JPY 1,000,000).
\end{itemize}

\section{Patents}
\hypertarget{sec:patents}{}
\noindent \textbf{Note}: \textit{These patents were filed during my time at Dalian University of Technology. In U.S. patents, my name appears as ``Deyun Lv", following pinyin conventions. In Chinese patents, my name appears as ``吕德运". My official academic name is ``Deyun Lyu".}

\subsection{U.S. Patents}
\begin{itemize}
    \item \textbf{[US-Pt3]} Weiqiang Kong, \textbf{Deyun Lyu}, Zhong Wei, Risheng Liu, Xin Fan and Zhongxuan Luo. ``\href{https://patents.google.com/patent/US20220174256A1/en}{Method For Infrared Small Target Detection Based On Depth Map In Complex Scene.}" US Patent, No. 12,108,022, Oct. 2024. (Granted, Serve as the primary contributor)
    \item \textbf{[US-Pt2]} Weiqiang Kong, \textbf{Deyun Lyu}, Zhong Wei, Risheng Liu, Xin Fan and Zhongxuan Luo. ``\href{https://patents.google.com/patent/US20220148213A1/en}{Method for Fully Automatically Detecting Chessboard Corner Points.}" US Patent, No. 12,094,152, Sept. 2024. (Granted, Serve as the primary contributor)
    \item \textbf{[US-Pt1]} Zhong Wei, \textbf{Deyun Lyu}, Weiqiang Kong, Risheng Liu, Xin Fan, Zhongxuan Luo and Shengquan Li. ``\href{https://patents.google.com/patent/US11900634B2/en}{Method for Adaptively Detecting Chessboard Sub-Pixel Level Corner Points.}" US Patent, No. 11,900,634, Feb. 2024. (Granted, Serve as the primary contributor)
\end{itemize}

\subsection{Chinese Patents}

\begin{itemize}
    \item \textbf{[CN-Pt6]} Weiqiang Kong, \textbf{Deyun Lyu}, Wei Zhong, Risheng Liu, Xin Fan, Zhongxuan Luo. ``\href{https://patents.google.com/patent/CN111243032B/en}{A Fully Automatic Checkerboard Corner Detection Method}", CN Patent, No. CN111243032B, May 2023. (Granted, Serve as the primary contributor)
    \item \textbf{[CN-Pt5]} Weiqiang Kong, \textbf{Deyun Lyu}, Wei Zhong, Risheng Liu, Xin Fan, Zhongxuan Luo. ``\href{https://patents.google.com/patent/CN111209877B/en}{An Infrared Small Target Detection Method Based on Depth Map in Complex Scenes}", CN Patent, No. CN111209877B, May 2023. (Granted, Serve as the primary contributor)
    \item \textbf{[CN-Pt4]} Wei Zhong, \textbf{Deyun Lyu}, Weiqiang Kong, Risheng Liu, Xin Fan, Zhongxuan Luo, Shengquan Li. ``\href{https://patents.google.com/patent/CN111260731B/en}{A Checkerboard Sub-pixel Level Corner Adaptive Detection Method}", CN Patent, No. CN111260731B, May 2023. (Granted, Serve as the primary contributor)
    \item \textbf{[CN-Pt3]} Wei Zhong, \textbf{Deyun Lyu}, Weiqiang Kong, Risheng Liu, Xin Fan, Zhongxuan Luo, Shengquan Li. ``\href{https://patents.google.com/patent/CN111291747B/en}{A Color Small Target Detection Method Based on Depth Map in Complex Scenes}", CN Patent, No. CN111291747B, June 2023. (Granted, Serve as the primary contributor)
    \item \textbf{[CN-Pt2]} \textbf{Deyun Lyu}, Wei Zhong, Weiqiang Kong, Risheng Liu, Xin Fan, Zhongxuan Luo. ``\href{https://patents.google.com/patent/CN111242991B/en}{Method for Quickly Registering Visible Light and Infrared Camera}", CN Patent, No. CN111242991B, Sept 2022. (Granted)
    \item \textbf{[CN-Pt1]} Wei Zhong, \textbf{Deyun Lyu}, Weiqiang Kong, Risheng Liu, Xin Fan, Zhongxuan Luo, Shengquan Li. ``\href{https://patents.google.com/patent/CN111243033B/en}{Method for Optimizing External Parameters of Binocular Camera}", CN Patent, No. CN111243033B, Sept 2022. (Granted, Serve as the primary contributor)
\end{itemize}

\section{Collaborative Research with Companies}
N/A (Not applicable for Assistant Professor position)

\section{Other Research Experience}

\subsection{Quality Assurance of End-to-End Autonomous Driving Systems}
{\customcventry{July 2025 -- Present}{Researcher}{National Institute of Informatics}{Tokyo, Japan}{}
{Collaboration with Toyota, \href{https://www.jari.or.jp/}{JARI} and \href{https://www.gaio.co.jp/}{GAIO}.
\begin{itemize}
    \item Demonstrated the potential of my prior research on end-to-end ADS.
    \item Developing new fault localisation and repair techniques for large-scale end-to-end controllers.
\end{itemize}
}}

\subsection{Testing of Multimodal Social Robots}
{\customcventry{June 2025 -- Present}{Researcher}{National Institute of Informatics}{Tokyo, Japan}{}
{Collaboration with University of Turin.
\begin{itemize}
  \item Ongoing development of a multimodal-aware metamorphic testing method of social robots (with \href{https://www.softbank.jp/robot/}{Pepper robot} as a case study).
\end{itemize}
}}

\subsection{eBPF-assisted Falsification of Temporal Logic Specifications under Noisy Environments}
{\customcventry{Apr. 2025 -- Present}{Researcher}{National Institute of Informatics}{Tokyo, Japan}{}
{Collaboration with Kyushu University and University of Bergamo.
\begin{itemize}
  \item Ongoing development of an eBPF-assisted falsification framework under noisy environments.
\end{itemize}
}}

\subsection{Benchmark Contribution to the \href{https://cps-vo.org/group/ARCH/FriendlyCompetition}{2025 ARCH Falsification Competition}}
{\customcventry{Apr. 2025 -- June. 2025}{Researcher}{National Institute of Informatics}{Tokyo, Japan}{}
{ARCH-COMP is a well-established, internationally recognized competition in CPS community.
\begin{itemize}
  \item Developed a novel benchmark generation method for falsification\hyperlink{pub:CPS-P3}{[CPS-P3]} and contributed five benchmarks to the 2025 ARCH Falsification Competition.
  \item Participated in the competition with our falsification tool \textit{ForeSee}\hyperlink{pub:CPS-P2}{[CPS-P2]}.
\end{itemize}
}}

\subsection{\texorpdfstring{\href{https://www.jst.go.jp/mirai/jp/program/super-smart/JPMJMI20B8.html}{Engineerable AI Project}}{}}
{\customcventry{Sept. 2023 -- Mar. 2024}{Visiting Researcher}{National Institute of Informatics}{Tokyo, Japan}{}
{University research project funding by JST Mirai project (Grant No. JPMJMI20B8).
\begin{itemize}
    \item Developed a fault localisation approach to localise faulty weights in DNN controllers.
    \item Proposed the first search-based repair approach for DNN controllers of AI-CPSs.
\end{itemize}
}}

\subsection{\texorpdfstring{\href{https://kaken.nii.ac.jp/en/grant/KAKENHI-PROJECT-23K28062/}{Testing, Analysis, and Repair of AI-Enabled Cyber-Physical Systems}}{}}
{\customcventry{Apr. 2023 -- Mar. 2025}{Research Assistant}{Kyushu University}{Fukuoka, Japan}{}
{University research project funded by JSPS (Grant No. 23H03372).
\begin{itemize}
    \item Designed an approach to localise faulty neurons in DNN controllers of AI-CPSs.
    \item Developed a data-driven method for automatically generating synthetic benchmarks for falsification.  
\end{itemize}
}}

\subsection{\texorpdfstring{\href{https://kaken.nii.ac.jp/en/grant/KAKENHI-PROJECT-21H04877/}{機械がバグを修正する時代―擬似オラクル生成・適用と自動バグ修正技術の深化}}{}}
{\customcventry{Jan. 2021 -- Aug. 2022}{Research Assistant}{Kyushu University}{Fukuoka, Japan}{}
{University research project funded by JSPS (Grant No. 21H04877).
\begin{itemize}
    \item Proposed eight time-aware coverage criteria and developed FalsifAI, the first coverage-guided falsification framework for AI-CPSs.
\end{itemize}
}}

\subsection{\texorpdfstring{\href{https://kaken.nii.ac.jp/en/grant/KAKENHI-PROJECT-20H04168/}{Comprehensive Analysis and Repairing Techniques for Stateful Deep Learning Systems}}{}}
{\customcventry{May 2021 -- Jan. 2022}{Research Assistant}{Kyushu University}{Fukuoka, Japan}{}
{University research project funded by JSPS (Grant No. 20H04168).
\begin{itemize}
    \item Established the first public benchmarks for AI-CPSs, enabling systematic evaluation and guiding testing research.
\end{itemize}
}}

\subsection{Traffic Accident Scene Reconstruction and Preprocessing Using TLS Point Clouds}
{\customcventry{Sept. 2019 -- June. 2020}{Research Assistant}{Dalian University of Technology}{Dalian, China}{}
{Collaboration with Guangzhou Traffic Police Department.
\begin{itemize}
    \item Developed a 3D point cloud-based system for traffic accident scene investigation and preprocessing.
    \item Implemented the system as a practical tool and deployed it in practice.
\end{itemize}
}}

\subsection{Vehicle-Mounted Visible–Infrared Stereo Vision Perception Unit}
{\customcventry{June 2018 -- Sept. 2019}{Research Assistant}{Dalian University of Technology}{Dalian, China}{}
{University–Industry Collaboration Project, resulting in multiple US and CN patents (See Section \hyperlink{sec:patents}{Patents}).
\begin{itemize}
    \item Designed automated sub-pixel chessboard corner detection methods and implemented them as practical calibration tools.
    \item Developed efficient external parameter optimization and r registration approaches for visible–infrared stereo vision perception systems.
    \item Proposed a depth map-based method for detecting small infrared/colour targets in complex scenes.
\end{itemize}
}}

\clearpage
\fancyhf{}
\renewcommand{\headrulewidth}{0.4pt}       % 页眉细线
\fancyhead[L]{\textsc{\color{gray}[5] Statement of Other Activities Related to Education and Research}} % 左页眉，小型大写+灰色
\fancyhead[R]{\textcolor{gray}{\thepage}}        % 右页眉显示页码

\section{Other Activities Related to Education and Research}

Beyond my individual research contributions, I view active participation in the academic and professional community as an essential part of my career. Since the beginning of my doctoral studies, I have consistently sought opportunities to contribute to the advancement of software engineering and AI-enabled CPSs through academic service, community engagement, and professional recognition. These activities complement my research and teaching, and they reflect my commitment to advancing both the science and practice of trustworthy intelligent systems.

\subsection{Academic Service}
I have served as a program committee member for major international conferences, including AAAI 2023-2025 (PC), ATVA 2025 (AE track), and I am the Web Chair of FM 2026. In addition, I regularly review for flagship journals such as ACM Transactions on Software Engineering and Methodology (TOSEM), as well as for leading conferences including ICSE, ASE, FSE, and ISSTA. Through these activities, I have contributed to upholding rigorous standards of research quality and fostering cross-disciplinary dialogue between software engineering, formal methods, and AI systems.

\subsection{Professional Awards}
My research contributions have been recognized by several professional awards. I received the IPSJ SIG-SE Excellent Research Award in 2025, one of the most prestigious recognitions in the Japanese software engineering community, which acknowledges the novelty and impact of my work on fault localization of AI-CPSs. In the same year, my dissertation was selected for the Recommended Doctoral Dissertation Award, recognizing the depth and originality of my doctoral research. I also received a GECCO 2024 Student Travel Grant, reflecting the international visibility of my research contributions. These honours not only validate the impact of my work but also motivate me to further strengthen my contributions to the community.

\subsection{Other Professional Contributions}
Beyond conferences and journals, I have contributed to the research community by developing publicly available benchmarks and tools. For instance, I co-developed the first industrial-level benchmark for AI-enabled CPSs, which has since been adopted by researchers worldwide. I also contributed five new benchmarks to the 2025 ARCH Falsification Competition, a long-standing international competition in the CPS community. These contributions have provided shared resources that enable fair evaluation and reproducibility, thereby shaping the research agenda of the community. In addition, I was appointed as a Global Talent Attraction Ambassador at Dalian University of Technology in 2024, further reflecting my engagement in promoting international research collaboration and education.

\subsection{Summary}
These activities, ranging from peer-review service and conference organization to benchmark development and international recognition, reflect my commitment to being an active academic citizen. Looking ahead, I aim to continue contributing to the community through program committees, journal editorial roles, and the creation of open resources, while also expanding collaborations with industry partners. I believe such efforts are indispensable to building a vibrant and trustworthy research ecosystem, and I am committed to advancing them alongside my own research and teaching.

\clearpage
\fancyhf{}
\renewcommand{\headrulewidth}{0.4pt}       % 页眉细线
\fancyhead[L]{\textsc{\color{gray}[6] Statement of Education and Research Plans in the Term of Five Years at JAIST}} % 左页眉，小型大写+灰色
\fancyhead[R]{\textcolor{gray}{\thepage}}        % 右页眉显示页码

\section{Education and Research Plan (Assistant Professor)}

% Describe concretely in 200-250 words your research plan by indicating the importance and novelty of the research and its value to society, and your specific plan for guiding and teaching students.

\textbf{For Research:} End-to-end autonomous driving (E2E AD) [\hyperlink{paper1}{1}-\hyperlink{paper3}{3}], a forefront instance of AI-CPSs, integrates perception, planning, and control into a single neural architecture. This black-box design obscures intermediate semantics (e.g., traffic lights), making it difficult to specify formal safety properties, monitor compliance, or localize faults. Moreover, the multimodal input space and large parameter scale result in a combinatorial explosion for testing \hyperlink{paper6}{[6]} and repair\hyperlink{paper9}{[9]}.

I will explore advanced neuro-symbolic approaches for the quality assurance of E2E AD\hyperlink{paper4}{[4},\hyperlink{paper5}{5]}. A key focus is developing methods to derive interpretable symbolic predicates from latent neural representations, enabling the exposure of meaningful events (e.g., pedestrian crossings). Building on this, I aim to establish multi-layer symbolic abstractions: neuron-level patterns, object-level semantics, and system-level driving behaviours. These abstractions will support the mining and synthesis of temporal logic specifications, reducing reliance on manually crafted rules. Moreover, I will investigate how symbolic constraints can be fed back into the neural controller to enable interpretable fault localisation \hyperlink{paper7}{[7},\hyperlink{paper8}{8]}, targeted repair, and runtime monitoring and shielding. Together, this neuro-symbolic loop integrates neural flexibility with formal rigour, advancing the reliability and trustworthiness of E2E AD and enhancing public acceptance of intelligent systems.

\textbf{For Education:} I will mentor graduate students to bridge theory and practice through hands-on projects with simulators and robotic platforms, emphasizing independent inquiry, rigorous experimentation, and scholarly writing. Regular seminars and collaborative projects will foster critical thinking across software engineering, machine learning, and formal methods, preparing them to contribute to trustworthy AI in society.


\vspace{6em} % 在 references 前留一行空白

{\footnotesize
\textbf{References} 
\begin{enumerate} [label={[\arabic*]}, leftmargin=*] 
  \item \hypertarget{paper1}{}Nikkei XTech. ``\href{https://xtech.nikkei.com/atcl/nxt/column/18/03080/012300002/}{自動運転に向けた次世代技術開発が加速}'' (in Japanese). Accessed Sep. 15, 2025.
  \item \hypertarget{paper2}{}B. Jiang et al. ``\href{https://openaccess.thecvf.com/content/ICCV2023/html/Jiang_VAD_Vectorized_Scene_Representation_for_Efficient_Autonomous_Driving_ICCV_2023_paper.html}{VAD: Vectorized Scene Representation for Efficient Autonomous Driving}", ICCV 2023.
  \item \hypertarget{paper3}{}X. Jia et al. ``\href{https://dl.acm.org/doi/10.5555/3737916.3737941}{Bench2Drive: Towards Multi-Ability Benchmarking of Closed-Loop End-to-End Autonomous Driving}", NeurIPS 2025.
  \item Wikipedia. \hypertarget{paper4}{}``\href{https://en.wikipedia.org/wiki/Neuro-symbolic_AI}{Neuro-symbolic AI}'' (in English). Accessed Sep. 15, 2025.
  \item J. Sun et al. \hypertarget{paper5}{}``\href{https://proceedings.mlr.press/v155/sun21a.html}{Neuro-Symbolic Program Search for Autonomous Driving Decision Module Design}", CoRL 2020. 
  \item Z. Zhang, D. Lyu et al. \hypertarget{paper6}{}``\href{https://ieeexplore.ieee.org/document/9844247}{FalsifAI: Falsification of AI-Enabled Hybrid Control Systems Guided by Time-Aware Coverage Criteria}", TSE 2023.
  \item \hypertarget{paper7}{}D. Lyu et al. ``\href{https://dl.acm.org/doi/10.1145/3705307}{SpectAcle: Fault Localisation of AI-Enabled CPS by Exploiting Sequences of DNN Controller Inferences}", TOSEM 2025.
  \item \hypertarget{paper8}{}D. Lyu et al. ``\href{https://www.sciencedirect.com/science/article/abs/pii/S0164121225001438}{Fault Localisation of AI-Enabled Cyber-Physical Systems by Exploiting Temporal Neuron Activation}", JSS 2025.
  \item \hypertarget{paper9}{}D. Lyu et al. ``\href{https://dl.acm.org/doi/10.1145/3638529.3654078}{Search-Based Repair of DNN Controllers of AI-Enabled CPS Guided by System-Level Specifications}", GECCO 2024.
 \end{enumerate}
}


\clearpage
\fancyhf{}
\renewcommand{\headrulewidth}{0.4pt}       % 页眉细线
\fancyhead[L]{\textsc{\color{gray}[7] Contact Information of Applicant}} % 左页眉，小型大写+灰色
\fancyhead[R]{\textcolor{gray}{\thepage}}        % 右页眉显示页码

\setlist[itemize]{leftmargin=0em,itemsep=0.15em,topsep=0.2em}
\section{Contact Information}
\begin{itemize}[label={}]
  \item \textbf{{Lyu Deyun}} \\
        Affiliation: Information Systems Architecture Science Research Division, National Institute of Informatics (NII), Tokyo, Japan \\
        Mailing Address: 2-1-2 Hitotsubashi, Chiyoda-ku, Tokyo 101-8430, Japan \\
        Phone Number: 080-2567-8231\\
        Email: lyudeyun@nii.ac.jp
\end{itemize}


\clearpage
\fancyhf{}
\renewcommand{\headrulewidth}{0.4pt}       % 页眉细线
\fancyhead[L]{\textsc{\color{gray}[8] Contact Information of Three References}} % 左页眉，小型大写+灰色
\fancyhead[R]{\textcolor{gray}{\thepage}}        % 右页眉显示页码

\section{Referees}
\setlist[itemize]{leftmargin=1.2em,itemsep=0.15em,topsep=0.2em}
\begin{itemize}
  \item \textbf{\href{https://researchmap.jp/f-ishikawa}{Prof. Fuyuki Ishikawa}} \\
        Affiliation: Information Systems Architecture Science Research Division, National Institute of Informatics (NII), Tokyo, Japan \\
        Mailing Address: 2-1-2 Hitotsubashi, Chiyoda-ku, Tokyo 101-8430, Japan \\
        Phone Number: 03-4212-2000 (ext. 2675)\\
        Email: f-ishikawa@nii.ac.jp

  \item \textbf{\href{https://researchmap.jp/read0190694}{Prof. Jianjun Zhao}} \\
        Affiliation: Faculty of Information Science and Electrical Engineering, Kyushu University, Fukuoka, Japan \\
        Mailing Address: 744 Motooka, Nishi-ku, Fukuoka 819-0395, Japan\\
        Phone Number: 092-802-3625\\
        Email: zhao@ait.kyushu-u.ac.jp

  \item \textbf{\href{https://faculty.dlut.edu.cn/kongweiqiang/en/index/828473/list/index.htm}{Prof. Weiqiang Kong}} \\
        Affiliation: School of Software Technology, Dalian University of Technology, Dalian, China \\
        Mailing Address: 321 Tuqiang Road, Jinpu New District, Dalian 116620, China\\
        Phone Number: +86-18840840913\\
        Email: wqkong@dlut.edu.cn     
\end{itemize}
\end{document}
